\documentclass[11pt]{article}
\usepackage[T1]{fontenc}
\usepackage[utf8]{inputenc}
\usepackage{lmodern}
\usepackage[margin=1in]{geometry}
\usepackage{amsmath,amssymb,mathtools}
\usepackage{tikz-cd}
\usepackage{microtype}
\usepackage[hidelinks]{hyperref}
\setlength{\parindent}{0pt}
\setlength{\parskip}{0.5em}
\allowdisplaybreaks
\newcommand{\C}{\mathbb C}
\newcommand{\R}{\mathbb R}
\newcommand{\PP}{\mathbb P}
\newcommand{\XX}{\mathcal X}
\newcommand{\OO}{\mathcal O}
\newcommand{\JJ}{\mathcal J}
\newcommand{\EE}{\mathcal E}
\newcommand{\UU}{\mathcal U}
\newcommand{\NN}{\mathcal N}
\newcommand{\id}{\mathrm{id}}
\newcommand{\bdel}{\bar\partial}
\DeclareMathOperator{\End}{End}
\DeclareMathOperator{\Lie}{Lie}
\DeclareMathOperator{\PGL}{PGL}
\title{Deformation of complex structures}
\author{}
\date{}
\begin{document}
\begin{center}
{\Large Deformation of complex structures}
\end{center}

% Source: handwritten PDF, page 1.
\paragraph{Motivation.}
Last time we had K\"ahler structure; today we do not.
What happens to Hodge structure when we perturb complex structures?

\paragraph{Motivating example: complex hypersurfaces.}
Let $X\subset\PP^n_{\C}$. Since $X$ is a hypersurface,
\[
 X=\{f=0\},\qquad \deg f=d.
\]
We want to deform $X$ to
\[
 X_t=\{f+tG=0\}\subset\PP^{n+1}_{\C},
 \qquad \deg G=\deg f=d,
\]
where $t\in\Delta\subset\C^n$, with $\Delta$ a polydisk.
Or, since $n=1$, $\Delta=\text{disk}$?

% Source: page 2.
\paragraph{Definition.}
\[
\begin{tikzcd}[column sep=large,row sep=large]
 \XX \arrow[r,hook]
     \arrow[dr,"\pi"']
 & \PP^{n+1}\times\Delta \arrow[d]\\
 & \Delta^{(n)}
\end{tikzcd}
\qquad \pi^{-1}(\{0\})=X.
\]
The embedding is holomorphic, and $\pi$ is a submersion.
$\Delta^{(n)}$ is a complex manifold.
$\XX$ should be our family of spaces.

\paragraph{What if I forget the ambient space $\PP^{n+1}_{\C}$?}
For example,
\[
 X=\{\underbrace{z_0z_1-z_2^2}_{f}=0\}\subset\PP^2_{\C},
 \qquad X_t=\{f+tG=0\},
\]
where $X\cong X_t\cong\PP^1_{\C}$.
So even though there is a deformation, the complex structure stays the
same. We want to also look at something different. This leads to the
definition:
\[
\begin{tikzcd}[column sep=large,row sep=large]
 X \arrow[r,hook] \arrow[d] & \XX \arrow[d,"\pi"]\\
 0 \arrow[r,phantom,"\in"] & \Delta
\end{tikzcd}
\qquad \pi\text{ a proper holomorphic submersion}.
\]
Ehresmann gives
\[
 \XX=X\times\Delta\qquad\text{(smooth trivializations)}.
\]
\[
 (X,J)\quad\leadsto\quad\{(X,J_t)\}\quad\text{for }t\in\Delta.
\]

% Source: page 3.
For $X\hookrightarrow\PP^{n+1}_{\C}$, get a short exact sequence of
vector bundles:
\[
 0\longrightarrow T_X\longrightarrow T_{\PP^{n+1}}|_X
 \longrightarrow
 \underbrace{T_{\PP^{n+1}}|_X/T_X}_{\NN_{X/\PP^{n+1}}}
 \longrightarrow0.
\]
Then get a long exact sequence of ?-cohomology:
\[
\begin{aligned}
 0&\longrightarrow H^0(X,T_X)
 \longrightarrow H^0(X,T_{\PP^{n+1}}|_X)\\
 &\longrightarrow H^0(X,\NN_{X/\PP^{n+1}})
 \xrightarrow{\delta}H^1(X,T_X).
\end{aligned}
\]
$\delta$: the Kodaira--Spencer map.
The cohomology theory is Dolbeault: \textit{see pic 2}.
% The right-hand side of this isomorphism is left blank in the source.
\[
 \underbrace{H^q(X,E)}_{(\mathrm{dR})}
 \underset{\text{(Dolbeault's theorem)}}{\cong}\quad
H^{p,q}_{\overline{\partial}}(X,E) = \frac{\text{closed}}{\text{exact}}\]

\paragraph{Main antagonist.}
\[
 H^0(X,\NN_{X/\PP^{n+1}})\xrightarrow{\delta}H^1(X,T_X).
\]
$H^0(X,\NN_{X/\PP^{n+1}})$: ``embedded first-order deformations of
$X\subset\PP^{n+1}_{\C}$.''

$H^1(X,T_X)$: ``abstract first-order deformations of $J$.''

$\ker\delta$: all first-order deformations of the complex structure
which are trivial.

% Source: page 4.
\paragraph{Question.}
Does every abstract deformation of $J$ arise from varying the embedding?
(The example from before is such an example of an element in the kernel.)

For a hypersurface in $\PP^{n+1}_{\C}$, let
\[
 S=\C[z_0,\ldots,z_{n+1}].
\]
See that it is graded:
\[
 S=\bigoplus_{k\geq0}S_k,
 \qquad S_k:\text{ homogeneous polynomials of degree }k,
 \qquad e_j=\left(\frac{z_j}{z_i}\right)^d e_i.
\]

\paragraph{Enter line bundles.}
Preview:
\[
 \C^{n+1}\ni\C z=\ell\qquad(z\neq0).
\]
For $d>0$,
\[
\begin{tikzcd}[row sep=large,column sep=small]
 \OO_{\PP^{n+1}}(d) \arrow[d]
 & \ni (\ell^{\vee})^{\otimes d}\\
 \PP^{n+1} & \ni[\ell]
\end{tikzcd}
\]
$(\ell^{\vee})^{\otimes d}$ corresponds to homogeneous degree-$d$
polynomials on $\ell$ (one-dimensional, so $\otimes=\operatorname{Sym}$).

From $S$, get local trivializations with transition map
\[
 e_j=\left(\frac{z_j}{z_i}\right)^d e_i
 \qquad\text{on }\{z_i\neq0\}.
\]
For $X\subset\PP^{n+1}_{\C}$,
\[
 \NN_{X/\PP^{n+1}}=\OO_X(d)\qquad(d>0).
\]
We are working on deforming hypersurfaces cut out by degree-$d$
homogeneous polynomials.

% Source: pages 4--5. The notation theta_X is retained from the notes.
\paragraph{Example: quartics in $\PP^3$ ($n=2$).}
\[
 H^0(X,\OO_X(4))\xrightarrow{\delta}H^1(X,\theta_X).
\]
\textit{Why $\theta_X$?}
\[
 H^0(X,\OO_X(4))\cong\C^{34},
\]
\[
 0\longrightarrow\ker\delta\hookrightarrow\C^{34}
 \xrightarrow{\delta}H^1(X,\theta_X)\longrightarrow\C\longrightarrow0.
\]
\[
 \ker\delta=\Lie(\PGL_4\C)\quad\text{(15-dimensional)},
 \qquad H^1(X,\theta_X)\quad\text{(20-dimensional)}.
\]
Details in Voisin II.6.

% Source: page 5.
\paragraph{Abstract deformations.}
Switch gears to abstract deformations, to see that they are parametrized
by $H^1(X,T_X)$.

$M=X$, but now with Riemannian structure to make sense of it; and it
won't matter which $g$.
\[
 \EE=\Gamma^{\infty}(X,\End_{\R}(TM)),
 \qquad
 \JJ(M)=\{J\in\EE\mid J^2=-\id\}.
\]
$\JJ(M)$: all possible complex structures; not a linear space.
The $C^{\infty}$ topology is given by seminorms
\[
 \|A\|_{C^r}=\max_{0\leq j\leq r}\sup_{x\in M}
 |\nabla^j A(x)|,
\]
where $M$ has a Levi-Civita connection $\nabla$ since $M$ has a
Riemannian metric $g$.

For $J\in\JJ(M)$,
\[
 T_{\C}M=T_J^{1,0}M\oplus T_J^{0,1}M,
\]
with projectors $\pi_J^{0,1}$ and $\pi_J^{1,0}$.
For another $J'\in\JJ(M)$, get
\[
 P_{J'}:=\pi_J^{0,1}|_{T_{J'}^{0,1}M}:
 T_{J'}^{0,1}M\longrightarrow T_J^{0,1}M.
\]

% Source: page 6. Type indices and directions are kept as written.
\[
 \UU_J=\{J'\in\JJ(M)\mid P_{J'}\text{ is a pointwise isomorphism}\},
 \qquad\UU_J\subset\JJ(M)\subset\EE.
\]
For $J'\in\UU_J$, $\xi_J\in A^{0,1}(M,T_J^{0,1}M)$, and
\[
 \xi_{J'}=\pi_J^{1,0}\circ P_{J'}^{-1}:
 T_J^{0,1}\longrightarrow T_J^{1,0}.
\]
Conversely, given
\[
 E=T_J^{1,0}\xrightarrow{\xi}F=T^{0,1},
\]
look at its graph
\[
 L_{\xi}=\{v+\xi(v)\mid v\in F\}.
\]
Define the complex conjugate map $\bar\xi:E\longrightarrow F$ and the
conjugate graph
\[
 \overline{L_{\xi}}=\{u+\bar\xi(u)\}.
\]
Then obtain $J'_{\xi}$ which does
\[
 J'_{\xi}u=-iu\quad\text{for }u\in L_{\xi},
\]
and then get
\[
 J_{\xi}u=iu\quad\text{for }u\in\overline{L_{\xi}},
 \qquad T_{\C}M=L_{\xi}\oplus\overline{L_{\xi}}.
\]
\paragraph{Problem.}
$J_{\xi}$ is not a complex structure, but an almost complex structure.
% The source writes +id here, although it writes -id earlier.
We merely have $J_{\xi}^2=\id$. Need to make it a complex structure.

Let $t=(t_1,\ldots,t_n)\in\Delta$, a polydisk. Look at
\[
 \xi_1(t)=\sum_{k=1}^{n}t^k\eta_k,
 \qquad \eta_1,\ldots,\eta_N:F\longrightarrow E,
\]
the first-order term of $\xi$. Get
\[
 \begin{aligned}
 \xi:T_0\Delta&\longrightarrow A^{0,1}(X,T_X),\\
 v=\left.\sum_kv^k\frac{\partial}{\partial t^k}\right|_{t=0}
 &\longmapsto\sum_kv^k\eta_k.
 \end{aligned}
\]

% Source: page 7.
\section*{Newlander--Nirenberg theorem}
If
\[
 \xi=\xi_{\bar i}^{a}\,d\bar z^{i}\times\partial_a,
 \qquad Z_i:=\partial_{\bar i}+\xi_{\bar i}^{a}\partial_a,
\]
then
\[
 [Z_i,Z_j]=C_{ij}^{a},
\]
so that we get
\[
 \bdel\xi+\frac12[\xi,\xi]=0
 \qquad\text{(Maurer--Cartan equation)}.
\]
So first-order deformations of complex structures correspond to
elements of $A^{0,1}(X,T_X)$ with $\bdel\eta=0$
(since, to first order, the bracket actually vanishes).

Now claim that
\[
 \xi:T_0\Delta\longrightarrow A^{0,1}(X,T_X)
\]
actually maps to
\[
 H^{0,1}_{\bdel}(X,T_X)=H^1(X,T_X)
 \qquad\text{(this is due to Kodaira--Spencer)}.
\]
I.e., if you

Since:
\begin{enumerate}
 \item closed with respect to $\bdel$;
 \item $\XX=X\times\Delta\longrightarrow\Delta$.
\end{enumerate}

\paragraph{Question.}
Given $\alpha\in H^1(X,T_X)$, can we obtain a family
$\XX\longrightarrow\Delta$?

Yes! But up to first order, provided $X$ has
\[
 H^2(X,T_X)=0.
\]
This is due to Kodaira--Nirenberg--Spencer.
\end{document}